\section*{Hoofdstuk \ref{ch:g11:prob} --- Kansrekening en toevalsvariabelen}

\begin{solution}{exo:g11:prob:1}
\omterm{def:g11:prob:model}{Uniform model} op $52$ kaarten. Harten: $\frac{13}{52} = \frac14$. Heren:
$\frac{4}{52} = \frac{1}{13}$. Harten of heer:
$\frac{13 + 4 - 1}{52} = \frac{16}{52} = \frac{4}{13}$ (de hartenheer telt
één keer mee). Geen van beide: $1 - \frac{4}{13} = \frac{9}{13}$.
\end{solution}

\begin{solution}{exo:g11:prob:2}
$S = 6$: de paren $(1,5), (2,4), (3,3), (4,2), (5,1)$, dus
$\P(S = 6) = \frac{5}{36}$.

$S \geq 10$: drie paren voor $10$, twee voor $11$, één voor $12$:
$\P(S \geq 10) = \frac{6}{36} = \frac16$.

$S$ even: de sommen $2, 4, 6, 8, 10, 12$ hebben samen
$1 + 3 + 5 + 5 + 3 + 1 = 18$ paren: \omterm{def:g10:proba:distribution}{kans} $\frac{18}{36} = \frac12$.
\end{solution}

\begin{solution}{exo:g11:prob:3}
De kansen tellen op tot $1$: $p = 1 - 0.3 - 0.2 - 0.4 = 0.1$. Dan is
\[
\E(X) = 0.3 \times (-1) + 0.2 \times 0 + 0.4 \times 2 + 0.1 \times 5
= -0.3 + 0.8 + 0.5 = 1 .
\]
\end{solution}

\begin{solution}{exo:g11:prob:4}
$\E(Y) = 2 \times 3 - 5 = 1$; $\V(Y) = 2^2 \times 4 = 16$;
$\sigma(Y) = 4$.
\end{solution}

\begin{solution}{exo:g11:prob:5}
De uitbetalingen zijn $0$, $2$ en $5$ euro met kansen $0.5$, $0.3$ en $0.2$;
na aftrek van de inzet van $2$ euro neemt $G$ de waarden $-2, 0, 3$ aan:
\begin{center}
\begin{tabular}{c|ccc}
$g$ & $-2$ & $0$ & $3$\\
\hline
$\P(G=g)$ & $0.5$ & $0.3$ & $0.2$
\end{tabular}
\end{center}
$\E(G) = -1 + 0 + 0.6 = -0.4$: de speler verliest gemiddeld $40$ cent per
spel --- ongunstig.
\end{solution}

\begin{solution}{exo:g11:prob:6}
Trekken in volgorde: $\P(X = 2) = \frac35 \times \frac24 = \frac{3}{10}$
(rood en dan rood); $\P(X = 0) = \frac25 \times \frac14 = \frac{1}{10}$
(blauw en dan blauw); bijgevolg is
$\P(X = 1) = 1 - \frac{3}{10} - \frac{1}{10} = \frac{6}{10}$. Dan is
\[
\E(X) = 0 \times \frac1{10} + 1 \times \frac6{10} + 2 \times \frac3{10}
= \frac{12}{10} = 1.2 ,
\]
in overeenstemming met de intuïtie: gemiddeld is $\frac35$ van de twee
trekkingen rood.
\end{solution}

\begin{solution}{exo:g11:prob:7}
$\E(G^2) = 0.5 \times 4 + 0.3 \times 0 + 0.2 \times 9 = 3.8$, dus geeft de
rekenformule $\V(G) = 3.8 - (-0.4)^2 = 3.8 - 0.16 = 3.64$ en
$\sigma(G) \approx 1.9$ euro.
\end{solution}

\begin{solution}{exo:g11:prob:8}
De winst is $80$ min de claim:
\begin{center}
\begin{tabular}{c|ccc}
$b$ & $80$ & $-920$ & $-4920$\\
\hline
$\P(B=b)$ & $0.94$ & $0.05$ & $0.01$
\end{tabular}
\end{center}
$\E(B) = 75.2 - 46 - 49.2 = -20$: tegen deze prijs \emph{verliest} de
maatschappij gemiddeld $20$ euro per polis. De premie is te laag voor de
gedekte risico's (ze zou boven $100$ euro moeten liggen voor een positieve
verwachte winst).
\end{solution}

\begin{solution}{exo:g11:prob:9}
De toevallige score is $+1$ met \omterm{def:g10:proba:distribution}{kans} $\frac14$ en $x$ met \omterm{def:g10:proba:distribution}{kans} $\frac34$:
\[
\E = \frac14 + \frac{3x}{4} = 0 \iff x = -\frac13 .
\]
Met de strafpunten $-\frac13$ levert willekeurig gokken gemiddeld niets op.
\end{solution}

\begin{solution}{exo:g11:prob:10}
Twee munten: ze vallen gelijk (KK of MM) met \omterm{def:g10:proba:distribution}{kans}
$\frac24 \times 2 = \frac12$. De winst van de speler is $+1$ of $-1$ met
gelijke kansen: \omterm{def:g11:prob:expectation}{verwachtingswaarde} $0$ voor beide spelers --- het spel is
eerlijk.

Drie munten: alle drie gelijk (KKK of MMM) met \omterm{def:g10:proba:distribution}{kans}
$\frac{2}{8} = \frac14$. De speler ontvangt $2$ met \omterm{def:g10:proba:distribution}{kans} $\frac14$ en betaalt
$1$ met \omterm{def:g10:proba:distribution}{kans} $\frac34$: $\E = \frac24 - \frac34 = -\frac14$. Het spel
bevoordeelt Ann met $25$ cent per ronde; het zou eerlijk zijn als ze $3$ euro
betaalde bij drie gelijke munten.
\end{solution}

\begin{solution}{exo:g11:prob:11}
\emph{1.} De organisator ontvangt $1000 \times 5 = 5000$ euro en keert in
totaal $2000 + 5 \times 200 + 20 \times 25 = 3500$ euro uit: de winst is de
\emph{constante} $1500$ (zodra alle loten verkocht zijn, is er geen toeval
meer). De winst $G$ van één speler (prijs min het lot van $5$ euro):
\begin{center}
\begin{tabular}{c|cccc}
$g$ & $1995$ & $195$ & $20$ & $-5$\\
\hline\rule{0pt}{11pt}%
$\P(G=g)$ & $\dfrac{1}{1000}$ & $\dfrac{5}{1000}$ & $\dfrac{20}{1000}$
& $\dfrac{974}{1000}$
\end{tabular}
\end{center}

\emph{2.}
\[
\E(G) = \frac{1995 + 5 \times 195 + 20 \times 20 - 974 \times 5}{1000}
= \frac{1995 + 975 + 400 - 4870}{1000} = -1.5 .
\]
Elke speler verwacht $1.50$ euro te verliezen; over $1000$ spelers is dat
$1500$ euro --- precies de winst van de organisator. De boeken kloppen: wat
de spelers gemiddeld verliezen, is precies wat de organisator int.
\end{solution}

\begin{solution}{pb:g11:prob:1}
\textbf{1.} $G$ neemt de waarden $-2, -1, 0, 1, 2, 3$ aan, elk met \omterm{def:g10:proba:distribution}{kans}
$\frac16$: $\E(G) = \frac{-2 - 1 + 0 + 1 + 2 + 3}{6} = \frac12$ euro.

\textbf{2.} $\E(G^2) = \frac{4 + 1 + 0 + 1 + 4 + 9}{6} = \frac{19}{6}$;
$V(G) = \frac{19}{6} - \frac14 = \frac{35}{12} \approx 2.92$;
$\sigma(G) \approx 1.71$.

\textbf{3.} $\E(2G - 1) = 2 \times \frac12 - 1 = 0$;
$V(2G - 1) = 4\,V(G) = \frac{35}{3}$.

\textbf{4.} $\E(\text{som}) = 7$: de eerlijke prijs is $7$ euro. Bij $8$ is
het verwachte verlies $1$ euro per spel.

\textbf{5.} Eerlijke premie: $0.05 \times 800 = 40$ euro. Bij $60$ euro:
verwachte winst $20$ euro per polis.

\textbf{6.} Verdelen in de verhouding $5:3$ beloont het \emph{verleden} ---
maar de pot behoort toe aan wie de \emph{toekomst} wint, en B kan die nog
winnen. Alles aan A geven behandelt een waarschijnlijke overwinning als
zeker --- de kleine \omterm{def:g10:proba:distribution}{kans} van B is iets waard.

\textbf{7.} Te spelen rondes: $3$; uitkomsten (A-winst/B-winst per ronde):
AAA, AAB, ABA, ABB, BAA, BAB, BBA, BBB. Speler B wint de reeks alleen door
alle drie te winnen: enkel BBB. A wint de reeks in $7$ van de $8$ gevallen:
eerlijke verdeling $64 \times \frac78 = 56$ pistolen voor A en $8$ voor B.

\textbf{8.} De extra, zinloze rondes spelen verandert niemands overwinning:
zodra A haar zesde winst heeft, zijn latere rondes plechtigheid. Maar de
horizon op precies $3$ rondes vastleggen maakt alle $8$ uitkomsten even
waarschijnlijk --- eerlijk tellen vraagt verhalen van gelijke lengte, en ze
opvullen kost niets.

\textbf{9.} Bij $5$--$5$: één beslissende ronde, $32$ elk. Bij $5$--$4$: A
krijgt $\frac{64 + 32}{2} = 48$. Bij $5$--$3$: de volgende ronde leidt naar
$64$ (A wint ze) of naar de stand $5$--$4$ die $48$ waard is: het aandeel van
A is $\frac{64 + 48}{2} = 56$ --- de $56 : 8$ van Fermat, van achteren naar
voren opnieuw afgeleid.

\textbf{10.} A heeft er $2$ nodig, B $3$: speel $4$ denkbeeldige rondes
($16$ uitkomsten). B wint de reeks bij $3$ of $4$ B-winsten:
$4 + 1 = 5$ uitkomsten. Verdeling: B: $64 \times \frac{5}{16} = 20$; A:
$44$.

\textbf{11.} De eerlijke huidige waarde van een onzekere toekomstige
uitbetaling is haar \omterm{def:g11:prob:expectation}{verwachtingswaarde} --- het met kansen gewogen \omterm{def:g11:stat:mean}{gemiddelde}
van wat komen kan. De verzekeringssector (vraag 5) en de financiële wereld
(het prijzen van onzekere uitbetalingen) zijn die ene zin, tot bedrijfstak
verheven; casino's ook, van de andere kant van de tafel.

\textbf{12.} De pot komt hoe dan ook ergens terecht: de twee aandelen zijn
\omterm{def:g11:prob:rv}{toevalsvariabelen} met som de constante $64$, dus geeft de lineariteit
$\E_A + \E_B = 64$ --- eerlijke aandelen putten de pot uit, bij elke stand.

\textbf{13.} Lineariteit: $\E(X + Y) = 3.5 + 3.5 = 7$, zonder tabel --- en
zelfs waar voor aan elkaar gelijmde dobbelstenen. \omterm{def:g11:prob:variance}{Variantie} van één
dobbelsteen: $\E(X^2) - 3.5^2 = \frac{91}{6} - \frac{49}{4} = \frac{35}{12}$;
voor onafhankelijke dobbelstenen is
$V(X + Y) = \frac{35}{12} + \frac{35}{12} = \frac{35}{6}$.

\textbf{14.} Elke vaste gast krijgt zijn eigen hoed terug met \omterm{def:g10:proba:distribution}{kans}
$\frac1n$ (zijn hoed is er één van de $n$). Optellen over de $n$ gasten
geeft $\E(X) = n \times \frac1n = 1$ --- gemiddeld één gelukkige gast, of
het feest nu twee hoeden telt of tweeduizend.

\textbf{15.} $n = 2$: de toewijzingen zijn ``allebei juist'' ($X = 2$) of
``verwisseld'' ($X = 0$): $\E = 1$. $n = 3$: de zes toewijzingen geven
$X = 3, 1, 1, 1, 0, 0$: $\E = \frac{6}{6} = 1$.

\textbf{16.} De berekening telde de \omterm{def:g11:prob:expectation}{verwachtingswaarden} van de $n$
afzonderlijke indicatoren op (``gast $i$ kreeg zijn hoed''), en de
lineariteit telt \omterm{def:g11:prob:expectation}{verwachtingswaarden} \emph{onvoorwaardelijk} op ---
afhankelijkheid zou een product of een \omterm{def:g11:prob:variance}{variantie} verwoesten, maar nooit een
som van \omterm{def:g11:prob:expectation}{verwachtingswaarden}. Dat is de superkracht van de lineariteit: er
worden geen vragen over onafhankelijkheid gesteld.

\textbf{17.} Beide winsten: $\E = 0.5 \times 10 - 5 = 0$ en
$0.005 \times 1000 - 5 = 0$. \omterm{def:g11:prob:variance}{Varianties} (van de winst): A:
$\E(G^2) = 0.5 \times 25 + 0.5 \times 25 = 25$; B:
$0.005 \times 995^2 + 0.995 \times 25 \approx 4\,975$: $\sigma_A = 5$ en
$\sigma_B \approx 70.5$. Mensen kopen B: wat ze kopen is niet de
\omterm{def:g11:prob:expectation}{verwachtingswaarde} (nul in beide gevallen) maar de \emph{spreiding} --- een
kleine \omterm{def:g10:proba:distribution}{kans} op een rechterstaart die een leven verandert.

\textbf{18.} De eerste munt bij worp $n$ heeft \omterm{def:g10:proba:distribution}{kans} $2^{-n}$ en betaalt
$2^n$ uit: elke $n$ draagt $2^{-n} \times 2^n = 1$ bij aan de
\omterm{def:g11:prob:expectation}{verwachtingswaarde}, die dus $1 + 1 + 1 + \dots$ is: oneindig. Toch betaalt
geen verstandig mens $1\,000$ om één keer te spelen: met \omterm{def:g10:proba:distribution}{kans}
$\frac{15}{16}$ is de uitbetaling hoogstens $16$. De \omterm{def:g11:prob:expectation}{verwachtingswaarde} vat
de lange reeks van \emph{herhaalbare} weddenschappen samen; een spel dat je
één keer in je leven speelt en dat een wilde staart heeft, prijs je niet met
haar \omterm{def:g11:stat:mean}{gemiddelde}.

\textbf{19.} \omterm{def:g11:prob:expectation}{Verwachtingswaarden}: verzekeren $-20$; onverzekerd blijven
$-0.01 \times 1000 = -10$. Gemiddeld wint de onverzekerde strategie --- voor
wie het verlies kan opvangen en de weddenschap vaak kan herhalen. Een gezin
dat de $-1\,000$ niet overleeft, zit niet in het middelen: het betaalt $10$
extra om de \omterm{def:g11:prob:variance}{variantie} te schrappen. Verzekeraars middelen; huishoudens kopen
zekerheid.

\textbf{20.} Geboren bij het verdelen van een pot tussen twee ongeduldige
gokkers; gesterkt door de lineariteit, die hoop optelt zonder de
afhankelijkheid om toestemming te vragen; blind voor het risico, dat de
\omterm{def:g11:prob:variance}{variantie} --- de loterijen, het Sint-Petersburgspel, de polis van het gezin
--- in de plaats moet meten; en volgend jaar gekroond door de wet van de
grote aantallen, die uitlegt waarom de \omterm{def:g11:stat:mean}{gemiddelden} uiteindelijk altijd
incasseren.
\end{solution}
